% chapter: wiki_antennas
\chapter{Antenna Fundamentals}\label{ch:wiki_antennas}

An antenna is a transducer that converts between guided electromagnetic waves (in a transmission line) and free-space radiation. Every antenna has a reciprocal receive and transmit behaviour described by the same parameters.

				

\section{Key Parameters}

				
\begin{itemize}

					\item \textbf{Gain (G)} — ratio of radiation intensity in the peak direction to that of an isotropic radiator, measured in dBi.

					\item \textbf{Directivity (D)} — same as gain but ignoring losses. G = η·D where η is radiation efficiency.

					\item \textbf{Beamwidth (HPBW)} — angular width of the main beam between half-power (−3 dB) points.

					\item \textbf{Bandwidth} — frequency range over which VSWR \textless{} 2 (or another specification).

					\item \textbf{Polarisation} — orientation of the E-field: linear, circular, or elliptical.

					\item \textbf{Radiation resistance} — equivalent resistance that would dissipate the same power as is radiated.

				
\end{itemize}

				

\section{Half-Wave Dipole}

				

\rffig{wiki_antennas_1.pdf}{Left: dipole E-plane figure-8 radiation pattern (gain ≈ 2.15 dBi). Right: rectangular patch antenna geometry.}

				

The reference antenna. Total length ≈ λ/2 (slightly less due to end effects). Gain = 2.15 dBi. Radiation resistance ≈ 73 Ω. Omnidirectional in the plane perpendicular to its axis.

\rffig{wiki_antennas_2.pdf}{Dipole radiation pattern: maximum at broadside (90°), nulls at 0° and 180°.}

				

\section{Patch (Microstrip) Antenna}

				

A rectangular conductor on a PCB substrate, resonant when the patch length L ≈ λ/(2√εreff). Gain ≈ 5–8 dBi, bandwidth 2–5\% (can be increased with thicker substrates). Widely used in GPS, WiFi, and cellular handsets due to its flat, PCB-compatible form factor.

				

\section{Array Factor}

				

When N identical elements are arranged in a line with spacing d and progressive phase shift δ, the combined pattern is the element pattern multiplied by the array factor AF(θ). The main beam can be steered electronically by changing δ without physically moving the array — the principle of phased array radar and 5G beamforming.
